\documentclass[preprint,12pt]{elsarticle}

\usepackage[utf8]{inputenc}
\usepackage[T1]{fontenc}
\usepackage{lmodern}

\usepackage{amsmath,amssymb,amsthm,mathtools}
\usepackage{enumitem}
\usepackage{booktabs}
\usepackage{array}
\usepackage{microtype}

\usepackage{hyperref}
\usepackage{xurl}
\usepackage[nameinlink,noabbrev]{cleveref}

\journal{Journal of Number Theory}

\hypersetup{
    colorlinks=true,
    linkcolor=blue,
    citecolor=blue,
    urlcolor=blue,
    pdftitle={An Explicit Evaluation of a Fibonacci Character Sum for
              Primes with Full Rank of Apparition},
    pdfauthor={Majid Ghandali},
    pdfsubject={Quadratic character sums of Fibonacci numbers under the
                full-rank hypothesis},
    pdfkeywords={Fibonacci numbers, rank of apparition, character sums,
                 primitive roots, finite fields}
}

\pdfstringdefDisableCommands{%
  \def\tnoteref#1{}%
  \def\corref#1{}%
  \def\F{F}%
  \def\Z{Z}%
  \def\Q{Q}%
  \def\R{R}%
  \def\N{N}%
  \def\QR{QR}%
  \def\NQR{NQR}%
}

\allowdisplaybreaks
\numberwithin{equation}{section}

\emergencystretch=1em
\Urlmuskip=0mu plus 1mu\relax
\hfuzz=1pt

\theoremstyle{plain}
\newtheorem{theorem}{Theorem}[section]
\newtheorem{lemma}[theorem]{Lemma}
\newtheorem{proposition}[theorem]{Proposition}
\newtheorem{conjecture}[theorem]{Conjecture}

\theoremstyle{definition}
\newtheorem{remark}[theorem]{Remark}

\crefname{lemma}{Lemma}{Lemmas}
\crefname{proposition}{Proposition}{Propositions}
\crefname{theorem}{Theorem}{Theorems}
\crefname{remark}{Remark}{Remarks}
\crefname{conjecture}{Conjecture}{Conjectures}

\renewcommand{\qedsymbol}{$\blacksquare$}

\DeclareMathOperator{\ord}{ord}

\newcommand{\Z}{\mathbb{Z}}
\newcommand{\Q}{\mathbb{Q}}
\newcommand{\R}{\mathbb{R}}
\newcommand{\F}{\mathbb{F}}
\newcommand{\N}{\mathbb{N}}
\newcommand{\Leg}[2]{\left(\frac{#1}{#2}\right)}
\newcommand{\eps}{\varepsilon}
\newcommand{\QR}{\mathrm{QR}}
\newcommand{\NQR}{\mathrm{NQR}}

\bibliographystyle{elsarticle-num}

\begin{document}

\begin{frontmatter}

\title{An Explicit Evaluation of a Fibonacci Character Sum for
Primes with Full Rank of Apparition\tnoteref{tn:preprint}}

\tnotetext[tn:preprint]{The preprint version of this article is
archived at Zenodo under DOI:
\href{https://doi.org/10.5281/zenodo.22803009}
{https://doi.org/10.5281/zenodo.22803009}.}

\author[inst1]{Majid Ghandali\corref{cor1}}
\cortext[cor1]{Corresponding author}
\ead{majid.ghandali@gmail.com}

\address[inst1]{Independent Researcher, Tehran, Iran}

\begin{abstract}
Let $p\ge 7$ be a prime, let $\chi_p$ denote the quadratic
character modulo $p$, and let $(F_n)$ be the Fibonacci sequence.
We evaluate the character sum
\[
S(p)=\sum_{n=1}^{p-1}\chi_p(F_n)
\]
under the hypothesis $\alpha(p)=p-1$, where $\alpha(p)$ denotes
the rank of apparition of $p$ in the Fibonacci sequence.

The proof identifies a structural mechanism specific to this
full-rank regime. The hypothesis forces $x^2-x-1$ to split over
$\F_p$, implies $p\equiv3\pmod4$, and makes the nonresidue root of
$x^2-x-1$ a primitive root of $\F_p^\times$. This yields a complete
multiplicative parametrization of the Fibonacci indices and reduces
the original sum to the single character value
\[
S(p)=\chi_p(r_- - r_+),
\]
where $r_-$ and $r_+$ are respectively the nonresidue and residue
roots of $x^2-x-1$. A cyclotomic argument involving fifth roots of
unity then gives
\[
S(p)=
\begin{cases}
+1,&p\equiv11\pmod{20},\\
-1,&p\equiv19\pmod{20}.
\end{cases}
\]
The condition $\alpha(p)=p-1$ may be viewed as a Fibonacci analogue
of an Artin-type primitivity condition. Independent computations
agree with the formula for all primes $p\le2\times10^6$ satisfying
$\alpha(p)=p-1$.
\end{abstract}

\begin{keyword}
Fibonacci numbers \sep Lucas sequences \sep character sums
\sep rank of apparition \sep primitive roots \sep fifth roots of unity

\MSC[2020] 11B39 \sep 11A07 \sep 11A15 \sep 11L40 \sep 11Y60
\end{keyword}

\hbadness=10000
\end{frontmatter}
\hbadness=1000

\section{Introduction}\label{sec:intro}
 
The arithmetic of the Fibonacci sequence modulo primes is organized
around two classical invariants: the \emph{Pisano period} $\pi(p)$
and the \emph{rank of apparition}
\[
\alpha(p)
=
\min\{m\ge1:\;p\mid F_m\}.
\]
Beginning with Lucas \cite{Lucas1878} and continuing through Wall
\cite{Wall1960}, Robinson \cite{Robinson1963}, and Vinson
\cite{Vinson1963}, a central theme has been the relationship between
these invariants and the multiplicative structure of finite fields.
For primes $p\ge7$, one has
\[
\alpha(p)\mid p-\chi_p(5).
\]
The extremal value $\alpha(p)=p-1$ can occur only in
the split case $\chi_p(5)=1$, which is equivalent to
$p\equiv\pm1\pmod5$. This extremal situation will be called the
\emph{full-rank regime} throughout the paper.
Rather than studying the period or rank themselves, we investigate
the quadratic character sum
\[
S(p)
=
\sum_{n=1}^{p-1}\chi_p(F_n)
\]
and ask whether the extremal condition $\alpha(p)=p-1$ imposes a
rigid structure on this sum.
 
At first sight there is little reason to expect such rigidity. Fibonacci congruences, primitive-root questions, and
quadratic-residuosity phenomena have been studied in several related
settings \cite{Williams1982,SunSun1992,Sun2019}. The principal
observation of the present paper is that the full-rank hypothesis
imposes substantially more structure than merely specifying the
value of $\alpha(p)$.
 
Indeed, the hypothesis produces a primitive-root parametrization of
the Fibonacci sequence. More precisely, the quadratic nonresidue root
of $x^2-x-1$ is shown to be a primitive root of $\F_p^\times$.
Once this is established, the Fibonacci index becomes a parametrization of the entire multiplicative group
$\F_p^\times$, reducing the original character sum to a single
quadratic character value.
 
Thus the proof follows the structural chain
\begin{equation}\label{eq:chain}
\boxed{
\begin{aligned}
\alpha(p) &= p-1 \\
&\Longrightarrow \text{primitive-root parametrization} \\
&\Longrightarrow \text{structural reduction} \\
&\Longrightarrow S(p)=\chi_p(r_- - r_+) \\
&\Longrightarrow S(p)=\pm1.
\end{aligned}
}
\end{equation}
 
From this perspective, the explicit value of the character sum is not
the starting point of the argument but its final consequence. The
central contribution is the structural mechanism that makes the
evaluation possible.
 
Our main theorem is as follows.
 
\begin{theorem}[Main Theorem]\label{thm:main}
Let $p\ge 7$ be a prime satisfying
\[
\alpha(p)=p-1.
\]
Then
\[
p\equiv11\pmod{20}
\qquad\text{or}\qquad
p\equiv19\pmod{20},
\]
and
\[
S(p)=
\begin{cases}
+1,& p\equiv11\pmod{20},\\
-1,& p\equiv19\pmod{20}.
\end{cases}
\]
\end{theorem}
 
The proof combines three ingredients. First, the full-rank condition
forces the nonresidue root of the Fibonacci polynomial to be a
primitive root modulo $p$ (\Cref{lem:full-rank-primitivity}).
Second, this primitive-root parametrization converts the Fibonacci
sequence into a complete enumeration of the multiplicative group,
causing the character sum to reduce to
$\chi_p(r_- - r_+)$ (\Cref{prop:structural-reduction}). Finally, a
cyclotomic discriminant criterion involving fifth roots of unity
determines this remaining character explicitly
(\Cref{lem:sign-determination}).
 
This viewpoint places the paper at the intersection of three
classical themes:
the arithmetic of Fibonacci sequences,
primitive-root phenomena of Artin type,
and quadratic character sums.
A detailed comparison with prior work, including the precise
relationship to the classical results of Shanks, De~Leon, and
Vinson, is given in the following subsection.
 
\subsection{Relation to the literature}
 
The arithmetic of Fibonacci and Lucas sequences modulo primes has
been studied from several complementary viewpoints; standard
references include \cite{Vajda1989,Koshy2001,Lemmermeyer2000,IrelandRosen1990}. Lucas \cite{Lucas1878} established basic
divisibility properties of Fibonacci numbers modulo primes, while
Wall \cite{Wall1960} initiated the systematic study of Pisano
periods. The relationship between the Pisano period $\pi(p)$ and the rank of
apparition $\alpha(p)$ was investigated further by Robinson
\cite{Robinson1963} and Vinson \cite{Vinson1963}.
 
Vinson \cite[Theorem~4(i)]{Vinson1963} proved that
\[
\frac{\pi(p)}{\alpha(p)}=1
\]
for primes $p\equiv11,19\pmod{20}$. Thus, in these residue classes,
the period and rank coincide. This result concerns the period--rank
quotient and does not evaluate the character sum considered here.
The present work concerns the distinct extremal condition
$\alpha(p)=p-1$. Its objective is not to determine the relationship
between $\pi(p)$ and $\alpha(p)$, but to evaluate the character sum
$S(p)$ under this rank-maximality hypothesis.
 
The primitive-root viewpoint adopted here is related to the notion of
a \emph{Fibonacci primitive root}, introduced by Shanks
\cite{Shanks1972} and developed further by Shanks--Taylor
\cite{ShanksTaylor1973} and De~Leon \cite{DeLeon1977}. De~Leon
\cite[Theorem~1]{DeLeon1977} proved that, for $p>5$, a Fibonacci
primitive root exists if and only if
\[
p\equiv1,9\pmod{10}
\quad\text{and}\quad
\pi(p)=p-1.
\]
 
In De~Leon's notation, $A(p)$ denotes the Pisano period. This
period-maximality condition is distinct from the rank-maximality
condition $\alpha(p)=p-1$ considered in the present paper. In general,
$\pi(p)=p-1$ and $\alpha(p)=p-1$ are not equivalent: for example,
$p=41$ has $\pi(41)=40$ and admits Fibonacci primitive
roots modulo $41$, whereas $\alpha(41)=20$.
A precise comparison between the full-rank condition $\alpha(p)=p-1$
and the Fibonacci-primitive-root condition, restricted to the split
case, is given in \Cref{prop:full-rank-fpr-equivalence} below.
 
The splitting behaviour of $x^2-x-1$ over finite fields and further
structural aspects of ranks and periods are considered in, among
other works, \cite{GuptaRockstrohSu2012,Renault2013,BallotElia2007,
SomerKrizek2015,FiebigMbirikaSpilker2025,JavaheriKrylov2020};
the limiting density of Fibonacci residues modulo prime powers has
recently been determined in \cite{BragmanRowland2025}.
Related congruence and quadratic-residuosity questions for Fibonacci
and Lucas sequences appear in \cite{Williams1982,SunSun1992,Sun2003,
BerndtEvansWilliams1998}.
 
The principal novelty of the present paper is therefore not the
classical theory of Fibonacci primitive roots or period-rank
relations. Rather, it is the structural reduction
\[
S(p)=\chi_p(r_- - r_+),
\]
obtained from the primitive-root parametrization forced by
$\alpha(p)=p-1$, together with the resulting explicit evaluation of
this character by a cyclotomic discriminant argument. We are not aware of any previous work establishing either this
reduction or the resulting formula for $S(p)$ in the full-rank regime.
 
The condition $\alpha(p)=p-1$ also places the problem in an
Artin-type setting. For the classical primitive-root problem, see
Hooley's conditional density theorem \cite{Hooley1967} and Moree's
survey \cite{Moree2012}. The density of primes admitting Fibonacci
primitive roots was studied conditionally by Sander
\cite{Sander1990}. These results concern related, but distinct,
primitive-root questions.
 
\subsection{Organisation}
 
\Cref{sec:preliminaries} fixes notation and records elementary
facts. \Cref{sec:structure} establishes the primitive-root structure
forced by the full-rank hypothesis. \Cref{sec:character-sums} and
\Cref{sec:structural-identity} derive the structural reduction of the
character sum. \Cref{sec:sign} determines the remaining sign, and
\Cref{sec:main-proof} proves \Cref{thm:main}. Computational checks
are reported in \Cref{sec:computations}, and
\Cref{sec:conclusion} discusses density and generalization questions.
 
\section{Preliminaries}\label{sec:preliminaries}
 
Throughout, $p$ denotes a prime with $p\ge7$, unless explicitly
stated otherwise. We use the standard
extension of the quadratic character to all of $\F_p$ by setting
$\chi_p(0)=0$, and recall that
\[
\chi_p(-1)=(-1)^{(p-1)/2}.
\]
We write
\[
\QR_p^\times
=
\{u\in\F_p^\times:\chi_p(u)=+1\},
\qquad
\NQR_p^\times
=
\{u\in\F_p^\times:\chi_p(u)=-1\}.
\]
 
When $p\equiv\pm1\pmod5$, the polynomial $x^2-x-1$ splits over
$\F_p$. Fix a square root $s\in\F_p$ of $5$, and write its two roots
as
\[
r_1=\frac{1+s}{2},
\qquad
r_2=\frac{1-s}{2}.
\]
Then
\[
r_1+r_2=1,
\qquad
r_1r_2=-1,
\qquad
r_1-r_2=s.
\]
The Binet formula holds in $\F_p$:
\begin{equation}\label{eq:binet}
F_m=\frac{r_1^m-r_2^m}{r_1-r_2}.
\end{equation}
 
After the quadratic-character analysis in
\Cref{lem:full-rank-primitivity}, in the full-rank regime we denote
the nonresidue root by $r_-$ and the residue root by $r_+$:
\[
\chi_p(r_-)=-1,
\qquad
\chi_p(r_+)=+1,
\]
\[
r_-+r_+=1,
\qquad
r_-r_+=-1.
\]
The Binet formula may then be written as
\[
F_m=\frac{r_-^m-r_+^m}{r_--r_+}.
\]
In \Cref{sec:sign}, the choice of $s$ is additionally made so that
$\chi_p(s)=+1$. This choice is independent of the labels $r_\pm$,
which are determined by quadratic character.
 
\subsection*{Elementary group-theoretic facts}
 
\begin{lemma}[Quadratic character in a cyclic group]
\label{lem:cyclic-character}
Let $g$ be a primitive root of $\F_p^\times$. Then
\[
\chi_p(g^k)=(-1)^k
\]
for every integer $k$.
\end{lemma}
 
\begin{proof}
The nonzero squares in $\F_p^\times$ are precisely the even powers
of $g$. Since $p-1$ is even, the parity of $k$ is well defined modulo
$p-1$.
\end{proof}
 
\begin{lemma}[Primitive-root reindexing]
\label{lem:primitive-reindexing}
Let $g$ be a primitive root of $\F_p^\times$. Then the map
\[
n\longmapsto g^n
\]
from $\{1,\dots,p-1\}$ to $\F_p^\times$ is a bijection. Consequently,
for every function $H:\F_p^\times\to A$, where $A$ is an abelian
group,
\[
\sum_{n=1}^{p-1}H(g^n)
=
\sum_{u\in\F_p^\times}H(u).
\]
Moreover,
\[
\chi_p(g^n)=(-1)^n
\qquad
(1\le n\le p-1).
\]
\end{lemma}
 
\begin{proof}
Since $g$ has order $p-1$, the powers
$g,g^2,\dots,g^{p-1}$ are exactly the elements of $\F_p^\times$,
each occurring once. The character formula follows from
\Cref{lem:cyclic-character}.
\end{proof}
 
\begin{lemma}[Fifth roots and the cyclotomic polynomial]
\label{lem:phi5-order}
Let $K$ be a field of characteristic different from $5$, and let
\[
\Phi_5(X)=X^4+X^3+X^2+X+1.
\]
For $z\in K$, one has $\Phi_5(z)=0$ if and only if
$z\in K^\times$, $z^5=1$, and $z\ne1$. In particular,
\[
\Phi_5(z)=0
\quad\Longrightarrow\quad
\ord(z)=5.
\]
\end{lemma}
 
\begin{proof}
The identity
\[
X^5-1=(X-1)\Phi_5(X)
\]
shows that $\Phi_5(z)=0$ implies $z^5=1$. Moreover,
$\Phi_5(1)=5\ne0$, so $z\ne1$. Conversely, if $z^5=1$ and
$z\ne1$, then
\[
(z-1)\Phi_5(z)=0
\]
and hence $\Phi_5(z)=0$.
\end{proof}
 
\section{Full-rank structure}\label{sec:structure}
 
\begin{lemma}[Rank--order identity]\label{lem:rank-order}
Let $r_1,r_2$ be the two roots of $x^2-x-1$ in $\F_{p^2}$, and put
\[
q=\frac{r_1}{r_2}.
\]
Then
\[
\alpha(p)=\ord_{\F_{p^2}^\times}(q).
\]
In particular, if $x^2-x-1$ splits over $\F_p$, then for each root
$r\in\F_p$,
\[
\alpha(p)=\ord_{\F_p^\times}(-r^2).
\]
\end{lemma}
 
\begin{proof}
Since $p\ne5$, the roots are distinct. Binet's formula in
$\F_{p^2}$ gives
\[
F_m\equiv0\pmod p
\quad\Longleftrightarrow\quad
r_1^m=r_2^m
\quad\Longleftrightarrow\quad
q^m=1.
\]
Thus
\[
\alpha(p)=\ord_{\F_{p^2}^\times}(q).
\]
 
Now suppose that $x^2-x-1$ splits over $\F_p$, and let $r$ be either
root. The other root is $-r^{-1}$, since the product of the roots is
$-1$. Thus the ratio $r_1/r_2$ is either $-r^2$ or its inverse,
depending on the ordering of the roots. These elements have the same
multiplicative order, so
\[
\alpha(p)=\ord_{\F_p^\times}(-r^2).
\]
\end{proof}
 
\begin{lemma}[Primitive-root mechanism in the full-rank regime]
\label{lem:full-rank-primitivity}
Let $p\ge7$ be a prime such that
\[
\alpha(p)=p-1.
\]
Then $x^2-x-1$ splits over $\F_p$, and
\[
p\equiv3\pmod4.
\]
Label its two roots $r_+,r_-\in\F_p^\times$ by their quadratic
characters:
\[
\chi_p(r_+)=+1,
\qquad
\chi_p(r_-)=-1.
\]
Then
\[
\ord_{\F_p^\times}(-r_-^2)=p-1
\]
and
\[
\ord_{\F_p^\times}(r_-)=p-1.
\]
In particular, $r_-$ is a primitive root of $\F_p^\times$.
\end{lemma}
 
\begin{proof}
Let $r_1,r_2$ be the two roots of $x^2-x-1$ in $\F_{p^2}$, and put
\[
q=\frac{r_1}{r_2}.
\]
 
Suppose first that $x^2-x-1$ does not split over $\F_p$. The
Frobenius automorphism of $\F_{p^2}$ interchanges its roots, so
\[
r_1^p=r_2,
\qquad
r_2^p=r_1.
\]
Consequently,
\[
q^p=\frac{r_1^p}{r_2^p}
=
\frac{r_2}{r_1}
=
q^{-1},
\]
and hence
\[
q^{p+1}=1.
\]
Thus $\ord(q)\mid p+1$. By \Cref{lem:rank-order},
\[
\alpha(p)\mid p+1.
\]
Since $\alpha(p)=p-1$, this gives $p-1\mid p+1$, hence
$p-1\mid2$, contradicting $p\ge7$. Therefore
$x^2-x-1$ splits over $\F_p$.
 
Let $r$ be either root in $\F_p$. By the split-case assertion of
\Cref{lem:rank-order},
\[
\ord(-r^2)=p-1.
\]
If $\chi_p(-1)=+1$, then $-r^2$ is a nonzero square and lies in the
subgroup of squares of order $(p-1)/2$, contradicting
$\ord(-r^2)=p-1$. Hence
\[
\chi_p(-1)=-1,
\]
or equivalently,
\[
p\equiv3\pmod4.
\]
 
Since $r_1r_2=-1$ and $\chi_p(-1)=-1$, the roots have opposite
quadratic characters. Hence precisely one root is a residue and the
other is a nonresidue; denote them by $r_+$ and $r_-$, respectively.
 
Applying \Cref{lem:rank-order} to $r_-$ gives
\[
\ord(-r_-^2)=\alpha(p)=p-1.
\]
Since $\chi_p(r_-)=-1$, Euler's criterion yields
\[
r_-^{(p-1)/2}=-1,
\]
so
\[
-1\in\langle r_-\rangle.
\]
Because $r_-^2\in\langle r_-\rangle$, it follows that
\[
-r_-^2\in\langle r_-\rangle.
\]
This element has order $p-1$, so $\langle r_-\rangle$ contains an
element of order $|\F_p^\times|=p-1$. Therefore
\[
\langle r_-\rangle=\F_p^\times,
\]
and hence
\[
\ord(r_-)=p-1.
\]
\end{proof}
 
\begin{remark}[Role of the full-rank hypothesis]
\label{rem:full-rank-role}
The essential consequence of
\Cref{lem:full-rank-primitivity} is that the substitution
\[
u=r_-^n
\]
is a bijection from $\{1,\dots,p-1\}$ onto $\F_p^\times$. This
allows the Fibonacci index to be replaced by a multiplicative
parameter in the proof of \Cref{prop:structural-reduction}.
\end{remark}
 
\begin{proposition}[Full rank and Fibonacci primitive roots]
\label{prop:full-rank-fpr-equivalence}
Let $p>5$ be a prime. Then, in the split case $\chi_p(5)=1$, the
following implications hold:
\[
\alpha(p)=p-1
\Longrightarrow
p\equiv3\pmod4 
\]
and $x^2-x-1$ has a Fibonacci primitive root modulo $p$.
Conversely, if $p\equiv3\pmod4$ and $x^2-x-1$ has a Fibonacci
primitive root modulo $p$, then
\[
\alpha(p)=p-1.
\]
\end{proposition}
\begin{proof}
The forward implication is the content of
\Cref{lem:full-rank-primitivity}: in the full-rank regime the
nonresidue root $r_-$ of $x^2-x-1$ satisfies
\[
\ord_{\F_p^\times}(r_-)=p-1,
\]
so $r_-$ is a Fibonacci primitive root.
 
Conversely, let $g$ be a Fibonacci primitive root. The other root of
$x^2-x-1$ is $-g^{-1}$, and the split-case identity of
\Cref{lem:rank-order} therefore gives
\[
\alpha(p)=\ord_{\F_p^\times}(-g^2).
\]
Since $p\equiv3\pmod4$ and $g$ is primitive,
\[
-1=g^{(p-1)/2},
\]
whence
\[
-g^2=g^{(p+3)/2}.
\]
Writing $p=4m+3$, we have $p-1=4m+2$ and
$(p+3)/2=2m+3$. Hence
\[
\gcd\!\left(p-1,\frac{p+3}{2}\right)
=
\gcd(4m+2,2m+3).
\]
Any common divisor divides
\[
(4m+2)-2(2m+3)=-4.
\]
Since $2m+3$ is odd, the common divisor is odd and divides $4$;
therefore it is equal to $1$. It follows that
\[
\ord_{\F_p^\times}(-g^2)=p-1,
\]
and consequently $\alpha(p)=p-1$.
\end{proof}
 
\begin{remark}[Relation to Fibonacci primitive roots]
\label{rem:primitive-root-literature}
\Cref{prop:full-rank-fpr-equivalence} makes the comparison with the
classical theory of Shanks \cite{Shanks1972} and De~Leon
\cite{DeLeon1977} completely explicit: in the split case, full rank
is equivalent to the existence of a Fibonacci primitive root together
with the congruence $p\equiv3\pmod4$. The direction used in the body
of the paper is the forward implication, which supplies the
primitive-root parametrization needed for the character-sum
evaluation.
\end{remark}
 
\begin{remark}[Order of the residue root]
\label{rem:residue-root-order}
In the full-rank regime, put
\[
n=\frac{p-1}{2}.
\]
Since $p\equiv3\pmod4$, the integer $n$ is odd. Since $r_-$ is
primitive,
\[
r_-^n=-1.
\]
Using $r_-r_+=-1$, we obtain
\[
r_+=-r_-^{-1}=r_-^{n-1}.
\]
Writing $p=4m+3$, so that $n=2m+1$, gives
\[
\ord(r_+)
=
\frac{p-1}{\gcd(p-1,n-1)}
=
\frac{4m+2}{\gcd(4m+2,2m)}
=
\frac{4m+2}{2}
=
\frac{p-1}{2}.
\]
\end{remark}
 
\section{Auxiliary character-sum identities}
\label{sec:character-sums}
 
The following elementary identities isolate the only character sums
required in the structural reduction.
 
\begin{lemma}[A quadratic character-sum identity]
\label{lem:quadratic-character-sum}
For every $a\in\F_p^\times$,
\[
\sum_{x\in\F_p}\chi_p(x^2-a)=-1.
\]
\end{lemma}
 
\begin{proof}
This is a standard quadratic-character-sum identity; see, for example,
\cite{IrelandRosen1990}. We include a proof for completeness.
 
Put
\[
T(a):=\sum_{x\in\F_p}\chi_p(x^2-a).
\]
Since
\[
\#\{x\in\F_p:x^2=t\}=1+\chi_p(t),
\]
we have
\[
T(a)
=
\sum_{t\in\F_p}
\bigl(1+\chi_p(t)\bigr)\chi_p(t-a)
=
\sum_{t\in\F_p}\chi_p\bigl(t(t-a)\bigr),
\]
because
\[
\sum_{t\in\F_p}\chi_p(t-a)=0.
\]
The substitution $t=aw$ shows that $T(a)$ is independent of
$a\ne0$. Since
\[
T(0)=\sum_{x\in\F_p}\chi_p(x^2)=p-1
\]
and
\[
\sum_{a\in\F_p}T(a)=0,
\]
the common value of $T(a)$ for $a\ne0$ is $-1$.
\end{proof}
 
\begin{lemma}[Negation interchanges residues and nonresidues]
\label{lem:negation-bijection}
If $p\equiv3\pmod4$, then the map
\[
u\longmapsto-u
\]
is a bijection
\[
\QR_p^\times\longrightarrow\NQR_p^\times.
\]
\end{lemma}
 
\begin{proof}
For $u\in\F_p^\times$,
\[
\chi_p(-u)=\chi_p(-1)\chi_p(u)=-\chi_p(u).
\]
The two sets have the same cardinality $(p-1)/2$.
\end{proof}
 
\begin{lemma}[Restricted quadratic character sums]
\label{lem:restricted-character-sums}
If $p\equiv3\pmod4$, then
\[
\sum_{u\in\QR_p^\times}\chi_p(u^2-1)=0
\]
and
\[
\sum_{u\in\NQR_p^\times}\chi_p(u^2+1)=-1.
\]
\end{lemma}
 
\begin{proof}
By \Cref{lem:quadratic-character-sum} with $a=1$, and isolating the
term $x=0$, we obtain
\[
\sum_{x\in\F_p^\times}\chi_p(x^2-1)
=
-1-\chi_p(-1)
=
0.
\]
Since $(-u)^2-1=u^2-1$ and negation maps
$\QR_p^\times$ bijectively onto $\NQR_p^\times$ by
\Cref{lem:negation-bijection},
\[
\sum_{u\in\QR_p^\times}\chi_p(u^2-1)
=
\sum_{v\in\NQR_p^\times}\chi_p(v^2-1).
\]
Thus each of these two sums is zero.
 
Similarly, \Cref{lem:quadratic-character-sum} with $a=-1$ gives
\[
\sum_{x\in\F_p^\times}\chi_p(x^2+1)
=
-1-\chi_p(1)
=
-2.
\]
Again, because $(-u)^2+1=u^2+1$ and negation interchanges the two
character classes,
\[
\sum_{u\in\NQR_p^\times}\chi_p(u^2+1)
=
\sum_{v\in\QR_p^\times}\chi_p(v^2+1).
\]
Therefore each of these two sums equals $-1$.
\end{proof}
 
\section{The structural reduction}\label{sec:structural-identity}
 
\begin{proposition}[Structural reduction of the character sum]
\label{prop:structural-reduction}
Let $p\ge7$ satisfy
\[
\alpha(p)=p-1.
\]
Then
\[
S(p)=\chi_p(r_--r_+),
\]
where $r_-$ and $r_+$ are respectively the nonresidue and residue
roots of $x^2-x-1$ in $\F_p$.
\end{proposition}
 
\begin{proof}
Put
\[
D=r_--r_+.
\]
Since the discriminant of $x^2-x-1$ is $5\ne0$ in $\F_p$, we have
$D\ne0$. Since
\[
r_+=-r_-^{-1},
\]
the Binet formula gives
\[
F_n
=
\frac{r_-^n-(-1)^n r_-^{-n}}{D}.
\]
 
By \Cref{lem:full-rank-primitivity}, $r_-$ is a primitive root of
$\F_p^\times$. Hence, by \Cref{lem:primitive-reindexing}, the map
\[
n\longmapsto u=r_-^n
\]
is a bijection from $\{1,\dots,p-1\}$ onto $\F_p^\times$, and
\[
\chi_p(u)=(-1)^n.
\]
Set
\[
\varepsilon=\chi_p(D).
\]
Then
\[
\chi_p(F_n)
=
\varepsilon
\chi_p\!\bigl(u-\chi_p(u)u^{-1}\bigr).
\]
 
If $u\in\QR_p^\times$, then $\chi_p(u)=+1$, and the contribution is
\[
\varepsilon\,\chi_p(u^2-1).
\]
If $u\in\NQR_p^\times$, then $\chi_p(u)=-1$, and the contribution is
\[
-\varepsilon\,\chi_p(u^2+1).
\]
Therefore, by \Cref{lem:restricted-character-sums},
\[
S(p)
=
\varepsilon
\left(
\sum_{u\in\QR_p^\times}\chi_p(u^2-1)
-
\sum_{u\in\NQR_p^\times}\chi_p(u^2+1)
\right)
=
\varepsilon(0-(-1)).
\]
Hence
\[
S(p)=\varepsilon=\chi_p(r_--r_+).
\qedhere
\]
\end{proof}
 
\begin{remark}
The term $n=p-1$, corresponding to $u=1$, contributes zero to the
sum through the factor $\chi_p(u^2-1)$ and therefore requires no
separate correction.
\end{remark}
 
\section{Evaluation of the remaining sign}\label{sec:sign}
 
Assume throughout this section that
\[
p\equiv3\pmod4
\qquad\text{and}\qquad
p\equiv\pm1\pmod5.
\]
Since $\chi_p(5)=1$, choose $s\in\F_p$ such that
\[
s^2=5
\qquad\text{and}\qquad
\chi_p(s)=+1.
\]
Such a choice is possible because the two square roots of $5$ have
opposite quadratic characters when $p\equiv3\pmod4$. Put
\[
\phi=\frac{1+s}{2},
\qquad
\psi=\frac{1-s}{2}.
\]
These are the two roots of $x^2-x-1$.
 
\begin{lemma}[Existence of fifth roots of unity]
\label{lem:fifth-root-existence}
There exists $z\in\F_p^\times$ with
\[
z^5=1
\qquad\text{and}\qquad
z\ne1
\]
if and only if
\[
p\equiv1\pmod5.
\]
\end{lemma}
 
\begin{proof}
The group $\F_p^\times$ is cyclic of order $p-1$. It contains an
element of order $5$ if and only if $5\mid p-1$.
\end{proof}
 
\begin{lemma}[Cyclotomic discriminant criterion]
\label{lem:cyclotomic-discriminant}
Let
\[
t=\frac{s-1}{2}
\qquad\text{and}\qquad
\Delta=t^2-4.
\]
Then $\Delta\ne0$ and
\[
\chi_p(\Delta)=
\begin{cases}
+1,&p\equiv1\pmod5,\\
-1,&p\equiv-1\pmod5.
\end{cases}
\]
\end{lemma}
 
\begin{proof}
Since $s^2=5$, one computes
\[
t^2+t-1=0
\]
and
\[
\Delta=-t-3.
\]
If $\Delta=0$, then $t=-3$, which together with
$t^2+t-1=0$ gives $5=0$ in $\F_p$, contradicting $p\ne5$.
Thus $\Delta\ne0$.
 
Suppose first that $\chi_p(\Delta)=+1$. Then the polynomial
\[
Z^2-tZ+1
\]
has a root $z\in\F_p$. Thus
\[
z+z^{-1}=t,
\]
and hence
\[
z^2+z^{-2}=t^2-2.
\]
Using $t^2+t-1=0$, we obtain
\[
(z^2+z^{-2})+(z+z^{-1})+1
=
(t^2-2)+t+1
=
0.
\]
Multiplying by $z^2$ yields
\[
\Phi_5(z)=0.
\]
By \Cref{lem:phi5-order},
\[
\ord(z)=5.
\]
Therefore $p\equiv1\pmod5$ by \Cref{lem:fifth-root-existence}.
 
Conversely, suppose that $p\equiv1\pmod5$. Choose
$z_0\in\F_p^\times$ of order $5$, and put
\[
\tau=z_0+z_0^{-1}.
\]
From $\Phi_5(z_0)=0$, one obtains
\[
\tau^2+\tau-1=0.
\]
The two roots of $X^2+X-1$ are $t$ and $-1-t$.
 
If $\tau=t$, then $z_0$ satisfies
\[
z_0^2-tz_0+1=0,
\]
so $\chi_p(\Delta)=+1$.
 
If $\tau=-1-t$, set $z_1=z_0^2$. Then $z_1$ also has order $5$, and
\[
z_1+z_1^{-1}
=
\tau^2-2
=
1-\tau-2
=
-1-\tau
=
t.
\]
Thus $z_1$ satisfies
\[
z_1^2-tz_1+1=0,
\]
and again $\chi_p(\Delta)=+1$.
 
We have proved that
\[
\chi_p(\Delta)=+1
\quad\Longleftrightarrow\quad
p\equiv1\pmod5.
\]
Since $\Delta\ne0$ and
$p\equiv\pm1\pmod5$, the remaining case gives
\[
\chi_p(\Delta)=-1
\]
when $p\equiv-1\pmod5$.
\end{proof}
 
\begin{lemma}[Sign determination]
\label{lem:sign-determination}
Let $r_-$ and $r_+$ denote the nonresidue and residue roots of
$x^2-x-1$, respectively. Then
\[
\chi_p(r_--r_+)=
\begin{cases}
+1,&p\equiv1\pmod5,\\
-1,&p\equiv-1\pmod5.
\end{cases}
\]
\end{lemma}
 
\begin{proof}
We have
\[
s\phi
=
s\frac{1+s}{2}
=
\frac{s+5}{2}.
\]
Since
\[
\Delta=-t-3
\qquad\text{and}\qquad
t=\frac{s-1}{2},
\]
it follows that
\[
\Delta
=
-\frac{s-1}{2}-3
=
-\frac{s+5}{2}
=
-s\phi.
\]
Taking quadratic characters and using
\[
\chi_p(s)=+1,
\qquad
\chi_p(-1)=-1,
\]
we obtain
\[
\chi_p(\Delta)=-\chi_p(\phi).
\]
 
Suppose that $p\equiv1\pmod5$. By
\Cref{lem:cyclotomic-discriminant},
\[
\chi_p(\Delta)=+1,
\]
so
\[
\chi_p(\phi)=-1.
\]
Thus
\[
r_-=\phi,
\qquad
r_+=\psi,
\]
and
\[
r_--r_+=\phi-\psi=s.
\]
Therefore
\[
\chi_p(r_--r_+)=+1.
\]
 
Now suppose that $p\equiv-1\pmod5$. Then
\[
\chi_p(\Delta)=-1,
\]
so
\[
\chi_p(\phi)=+1.
\]
Thus
\[
r_+=\phi,
\qquad
r_-=\psi,
\]
and
\[
r_--r_+=\psi-\phi=-s.
\]
Therefore
\[
\chi_p(r_--r_+)
=
\chi_p(-1)\chi_p(s)
=
-1.
\]
\end{proof}
 
\section{Proof of the main theorem}\label{sec:main-proof}
 
\begin{proof}[Proof of \Cref{thm:main}]
By \Cref{lem:full-rank-primitivity}, the polynomial $x^2-x-1$
splits over $\F_p$ and
\[
p\equiv3\pmod4.
\]
Splitting is equivalent to $\chi_p(5)=+1$, hence to
\[
p\equiv\pm1\pmod5.
\]
Combining these congruences gives
\[
p\equiv11\pmod{20}
\qquad\text{or}\qquad
p\equiv19\pmod{20}.
\]
 
By \Cref{prop:structural-reduction},
\[
S(p)=\chi_p(r_--r_+).
\]
By \Cref{lem:sign-determination}, this is $+1$ when
$p\equiv1\pmod5$, equivalently when $p\equiv11\pmod{20}$, and it is
$-1$ when $p\equiv-1\pmod5$, equivalently when
$p\equiv19\pmod{20}$.
\end{proof}
 
\section{Computational checks}\label{sec:computations}
 
The conclusion of \Cref{thm:main} was checked for every prime
\[
p\le2\times10^6
\]
satisfying
\[
\alpha(p)=p-1.
\]
 
\begin{table}[ht]
\centering
\caption{Computational check of \Cref{thm:main} for all primes
$p\le2\times10^6$ satisfying $\alpha(p)=p-1$.}
\label{tab:verification}
\small
\begin{tabular}{@{}ccccc@{}}
\toprule
$p\bmod 20$ & predicted $S(p)$ & count & matches & status\\
\midrule
$11$ & $+1$ & $11{,}755$ & $11{,}755$ & PASS\\
$19$ & $-1$ & $14{,}652$ & $14{,}652$ & PASS\\
\midrule
total & & $26{,}407$ & $26{,}407$ & PASS\\
\bottomrule
\end{tabular}
\end{table}
 
The verification is non-circular: both $\alpha(p)$ and $S(p)$
were computed directly from their definitions, and neither
calculation uses the closed formula proved in \Cref{thm:main}.
In particular, $\alpha(p)$ was obtained by iterating the Fibonacci
recurrence modulo $p$ and recording the least index $n\ge1$ for which
$F_n\equiv0\pmod p$, while $S(p)$ was obtained by summing
$\chi_p(F_n)$ for $1\le n\le p-1$. The subsequent consistency checks
compare these
direct computations with the structural predictions: (i)~the
quadratic characters of both roots of $x^2-x-1$ in $\F_p$ against
\Cref{lem:full-rank-primitivity}; (ii)~the multiplicative orders of
$-r_-^2$ and $r_-$ against
\Cref{lem:rank-order,lem:full-rank-primitivity}; and (iii)~the sign
of $\chi_p(r_--r_+)$ against the discriminant criterion of
\Cref{lem:sign-determination}. These checks provide independent
computational support for both the predicted value of $S(p)$ and
the intermediate identities used in the proof.
 
Among the $148{,}933$ primes $p\le2\times10^6$, exactly $26{,}407$
satisfy $\alpha(p)=p-1$. Among the $74{,}461$ split primes
\[
p\equiv\pm1\pmod5,
\]
the observed full-rank proportion is
\[
\frac{26{,}407}{74{,}461}\approx0.35459
\qquad(35.46\%).
\]
For comparison, the first-order Artin-type value suggested by the
classical classwise Fibonacci-primitive-root densities
(\Cref{rem:density-heuristic}) is
\[
C_F^{\mathrm{heur}}=\frac{18A}{19}\approx0.35427.
\]
The observed ratio differs from this heuristic value by about
$3.1\times10^{-4}$ in absolute terms (relative difference below
$0.1\%$). This comparison is empirical only and is not used in the
proof of \Cref{thm:main}. The computation was performed in Python.
Source code, datasets, and verification logs are available in the
reproducibility archive described in \Cref{app:software}.
 
\begin{remark}[Independent computational cross-check]
\label{rem:independent-cross-check}
A separate Python implementation, independent of the main codebase,
recomputed $\alpha(p)$ and $S(p)$ for every prime $p\le20{,}000$ and
for a random sample of $150$ primes satisfying
\[
20{,}000<p\le5\times10^5.
\]
The independent recomputation agreed with the main computation in
every case. Moreover, the proportion
\[
\frac{7{,}379}{20{,}751}\approx0.3556
\]
among split primes at the bound $5\times10^5$, as reported by the
main codebase, is consistent with the value $0.35459$ obtained at the
bound $2\times10^6$ and with the first-order heuristic
$C_F^{\mathrm{heur}}=18A/19\approx0.35427$ of
\Cref{rem:density-heuristic}.
\end{remark}
 
\section{Concluding remarks}\label{sec:conclusion}
 
We have obtained an explicit evaluation of the Fibonacci character
sum
\[
S(p)=\sum_{n=1}^{p-1}\chi_p(F_n)
\]
in the full-rank regime
\[
\alpha(p)=p-1.
\]
The proof shows that the sum is governed by the structural reduction
\[
S(p)=\chi_p(r_--r_+),
\]
which in turn arises from the primitive-root parametrization forced by
the full-rank condition. The remaining sign is determined by
$p\bmod 20$.
 
The split-case identity in \Cref{lem:rank-order} identifies
$\alpha(p)=p-1$ with the primitivity of
\[
-r^2\in\F_p^\times
\]
for either root $r$ of $x^2-x-1$. Thus the full-rank condition belongs
naturally to an Artin-type family of primitive-root problems.
In the split setting, \Cref{prop:full-rank-fpr-equivalence} further
identifies the full-rank condition with the Fibonacci-primitive-root
condition restricted to the two residue classes
$p\equiv11,19\pmod{20}$. This gives the conjectural density constant
\[
C_F^{\mathrm{heur}}=\frac{18A}{19}
\]
a classical Artin-type interpretation through the classwise
Fibonacci-primitive-root densities of Shanks~\cite{Shanks1972}.
 
\subsection*{A density question}
 
Among primes $p\le2\times10^6$ in the split classes
\[
p\equiv\pm1\pmod5,
\]
the observed full-rank proportion is
\[
\frac{26{,}407}{74{,}461}\approx0.35459
\]
(see the numerical verification above). By \Cref{thm:main}, such
primes can occur only in the two classes
\[
p\equiv11,19\pmod{20},
\]
among the four split residue classes
\[
p\equiv1,9,11,19\pmod{20}.
\]
This motivates the following conjecture.
 
\begin{conjecture}\label{conj:density}
Let
\[
N_F(x)
=
\#\{p\le x:\alpha(p)=p-1\}
\]
and
\[
N_{\mathrm{split}}(x)
=
\#\{p\le x:p\equiv\pm1\pmod5\}.
\]
There exists a constant $C_F$ with
\[
0<C_F\le\frac12
\]
such that
\[
N_F(x)\sim C_F\,N_{\mathrm{split}}(x)
\qquad
(x\to\infty).
\]
Equivalently, the natural density of full-rank primes among all primes
equals $C_F/2$.
\end{conjecture}
 
The upper bound
\[
C_F\le\frac12
\]
is forced by \Cref{thm:main}. By the split-case identity in
\Cref{lem:rank-order}, \Cref{conj:density} is a primitive-root
statement for the reduction of the algebraic unit
\[
-\phi^2,
\qquad
\phi=\frac{1+\sqrt5}{2},
\]
at primes split in $\Q(\sqrt5)$.
The primitive-root interpretation of
\Cref{prop:full-rank-fpr-equivalence} gives a classical route to a
first-order heuristic for the constant $C_F$. In the four split
classes
\[
p\equiv1,9,11,19\pmod{20},
\]
the full-rank condition can occur only in the two classes
\[
p\equiv11,19\pmod{20}
\]
by \Cref{thm:main}. The relevant conditional densities for Fibonacci
primitive roots in these two classes were computed class-by-class by
Shanks~\cite{Shanks1972}; see also the later conditional theorem of
Sander~\cite{Sander1990}. Thus the heuristic value of $C_F$ can be
obtained directly from the classical Fibonacci-primitive-root
densities, after restricting to the two classes allowed by the
full-rank theorem.
 
\begin{remark}[Heuristic value of $C_F$]
\label{rem:density-heuristic}
The constant $C_F$ is naturally related to the classical
Fibonacci-primitive-root densities. Shanks~\cite{Shanks1972} gives the
following first-order densities, conditional on the indicated residue
class modulo $20$:
\[
\begin{array}{c|cccc}
p\bmod20 & 1 & 9 & 11 & 19\\
\hline
\delta_{\mathrm{FPR}}(p\bmod20)
& \dfrac{16A}{19}
& \dfrac{20A}{19}
& \dfrac{32A}{19}
& \dfrac{40A}{19},
\end{array}
\]
where
\[
A=\prod_q\Bigl(1-\frac1{q(q-1)}\Bigr)\approx0.3739558136
\]
is Artin's constant. The remaining four residue classes modulo $20$
have zero density for Fibonacci primitive roots. Shanks's averaging
over all eight residue classes yields the classical constant
\[
\frac{27A}{38}
\]
for the density of primes admitting a Fibonacci primitive root.
 
For the present problem, however, only the classes
$p\equiv11,19\pmod{20}$ can contribute, by \Cref{thm:main}. Among the
four split classes
\[
p\equiv1,9,11,19\pmod{20},
\]
these two classes have asymptotic relative weight $1/2$. Conditional
on either of these two classes, the corresponding first-order
Fibonacci-primitive-root densities are
\[
\frac{32A}{19}
\quad\text{and}\quad
\frac{40A}{19}.
\]
Consequently, the resulting first-order conditional value in
\Cref{conj:density} is
\[
\begin{aligned}
C_F^{\mathrm{heur}}
&=
\frac12
\Biggl(
\frac12
\Biggl(
\frac{32A}{19}+\frac{40A}{19}
\Biggr)
\Biggr)
=
\frac{18A}{19}
\approx0.354274.
\end{aligned}
\]
Thus
\[
\boxed{C_F^{\mathrm{heur}}=\dfrac{18A}{19}}.
\]
This value is recorded here as a conjectural first-order constant,
not as a proved identity. The relevant analytic problem is an
Artin-type primitive-root problem for the algebraic setting arising
from the split Fibonacci polynomial. Sander~\cite{Sander1990}
established the classical Fibonacci-primitive-root density under
appropriate generalized Riemann hypotheses; a corresponding
derivation for the present conditional formulation is not developed
here.
\end{remark}
 
It would be natural to ask whether a Hooley-type analysis, under
suitable GRH assumptions, yields a rigorous justification of the
value $18A/19$. This question lies beyond the scope of the present
paper.
 
\subsection*{Further directions}
 
The same primitive-root parametrization remains available for Lucas-type sums
\[
L_n=r_-^n+r_+^n.
\]
It is therefore natural to investigate whether an analogue of the
reindexing argument in \Cref{prop:structural-reduction} yields
explicit evaluations for character sums attached to Lucas sequences;
for recent work on the arithmetic of Lucas sequences and the integers
they represent, see \cite{HajduTijdeman2025}.
This extension is not pursued here.
 
More generally, the present result suggests that extremal rank
conditions may provide a mechanism for evaluating character sums
attached to linear recurrence sequences. Extending this approach to
higher-order characters and to more general Lucas sequences is a
natural direction for future work.
 
\appendix
 
\section{Illustrative examples}\label{app:examples}
 
\begin{table}[ht]
\centering
\caption{Illustrative examples of \Cref{thm:main} for selected
primes $p$ satisfying $\alpha(p)=p-1$.}
\label{tab:examples}
\small
\begin{tabular}{@{}cccccc@{}}
\toprule
$p$ & $\alpha(p)$ & $p\bmod 20$ & $r_-$ & $r_+$ & $r_--r_+$\\
\midrule
$11$ & $10$ & $11$ & $8$ & $4$ & $4=2^2$ (QR)\\
$19$ & $18$ & $19$ & $15$ & $5$ & $10$ (NQR)\\
$31$ & $30$ & $11$ & $13$ & $19$ & $25=5^2\pmod{31}$ (QR)\\
$59$ & $58$ & $19$ & $34$ & $26$ & $8$ (NQR)\\
$79$ & $78$ & $19$ & $30$ & $50$ & $59$ (NQR)\\
\bottomrule
\end{tabular}
\end{table}
 
In each case, the quadratic character of $r_--r_+$ agrees with
\Cref{thm:main}: the value is $+1$ in the class $11\pmod{20}$ and
$-1$ in the class $19\pmod{20}$. Accordingly,
\begin{gather*}
S(11)=+1,\qquad S(19)=-1,\qquad S(31)=+1, \\
S(59)=-1,\qquad S(79)=-1.
\end{gather*}
 
\section{Reproducibility archive}\label{app:software}
 
The source code, datasets, and verification logs used for the
computations are available in the accompanying GitHub repository and
archived at~\cite{Ghandali2026}. The archive includes the principal
implementation, \texttt{fibchar\_v1\_0\_0.py}, together with the data
files required to reproduce the computational summary in
\Cref{tab:verification}.
 
\bibliography{references_full_rank}
 
\end{document}